\clearpage
\section{Variational crossing area and a uniform read action}
\label{sec:m1-interface-compactness}

Agreement on fixed smooth sections does not ensure that a crossing statistic
selects macroscopic interfaces. A sparse read family can retain the leading
planar area coefficient while admitting microscopic partitions whose entropy
lies below the continuum isoperimetric value. An explicit addition of shorter
reads restores compactness, the full perimeter limit and the limiting minimum
shape. The same scale threshold removes an extra sector of the uniform
read-action spectrum. These are statements about declared mathematical read
graphs, coordinate volumes and observables; no physical horizon or source
identification is an antecedent or a conclusion.

\paragraph{Finite definition.}
Set $L=1$ initially. Use $q^3$ equal cells of side $h=1/q$, either on the
unit torus or in the unit cube. In the latter case all occupied cells must
lie in a fixed interior cube $C=[\eta,1-\eta]^3$, $0<\eta<1/2$; reads are
clipped to the actual site population. For sufficiently small read radius,
the cut count then equals its zero-extended count without any outside record.
A binary site partition $u$ has unordered crossing count and normalized value
\[
 X_V(u)=\frac12\sum_x\sum_{v\in V}|u(x+v)-u(x)|,
 \qquad \mathcal E_V(u)=\frac{4X_V(u)}{\pi q^4}.
\]
Every unordered pair has one nat, and its two layer directions are counted
only once here. The physical area conversion is $L^2\mathcal E_V$.

For $n\ge2$ put
\[
 q=2^{8n},\quad m=2^{4n},\quad K=2^{3n},\quad R=mK,
 \qquad \epsilon=m/q,\quad a=R/q.
\]
The stencil $T$ consists of every $mv$, $v\in\mathbb Z^3$, $\|v\|\le K$,
and all signed dyadic axis steps below $m$, with zero included once. Its
previous-layer tick is $a/c$. The identities $q=m^2$, $mK^4=q^2$ and
$a^4=q\epsilon^3$ fix its coarse entropy as
\[
 \mathcal E_{\rm coarse}(u_h)
 =\frac{2}{\pi q}\sum_{\|v\|\le K}
       \|u_h(\cdot+\epsilon v)-u_h\|_{L^1},
\]
where $u_h$ is the piecewise-constant cell indicator. The ball flux
$\int_{\|z\|\le a}|z\cdot\nu|\,dz=\pi a^4/2$ gives coefficient one
for every fixed finite-perimeter interface.

\paragraph{A microscopic partition with the wrong entropy.}
On the torus occupy precisely the sites with $x_1\bmod m<m/2$. The volume
is exactly one half. Coarse reads preserve the residue label. A dyadic
$x$-step $s\le m/2$ changes a fraction $2s/m$ of labels, so
\[
 X_T=\frac{2q^3(m-1)}m,\qquad
 \mathcal E_T=\frac{8(m-1)}{\pi m q}\longrightarrow0.
\]
These indicators converge weakly to $1/2$ and have no strongly $L^1$
convergent subsequence. Thus bounded crossing entropy does not imply a
macroscopic region, and the finite half-volume minima tend to zero.

There is also a clipped-box violation of the expected minimum. Let a ball
$F\Subset C$ have volume $2v$, and retain only the active residue half in
$F$. Coarse displacements preserve that residue. Averaging it over its
$\epsilon$-period changes each crossing slab by at most its boundary area
times $O(\epsilon)$; the normalized error is $O(\epsilon/a)$.
The cell error is $O(h/a)$ and the dyadic contribution is
$O(\log q/q)$. Consequently
\[
 |u_h|\to v,\qquad
 \mathcal E_T(u_h)\to\tfrac12\operatorname{Per}(F)
 =2^{-1/3}(36\pi)^{1/3}v^{2/3}.
\]
This is strictly below the Euclidean isoperimetric value at volume $v$.
The count error is $O(mq^2)$ cells. Correcting it to $\operatorname{round}(vq^3)$
changes the normalized cut by at most $O(mD_T/q^2)=O(K^{-1})+o(1)$,
so the exact finite volume constraint retains the counterexample.

Connecting the occupied slabs by a one-cell-wide axial path changes only
$O(q)$ cells and adds entropy $O(D_T/q^3)\to0$. This connected sequence
has weak limit $\tfrac12\mathbf1_F$ and no strong subsequence.
Connectivity alone therefore does not restore compactness. This connected
control uses volume convergence; the exact-volume minimum above imposes
no separate connectivity constraint. Fixed-section area agreement is
consistent with all these failures, because it did not quantify over
microscopically varying cuts.

\paragraph{An explicit sparse repair.}
Set $U=T\cup B_r$, where $B_r=\{w\in\mathbb Z^3:\|w\|\le r\}$ and
$r\le m/2$. All read and entropy weights remain equal. The relevant window is
\[
 \frac{r^4}{mq}\longrightarrow\infty,
 \qquad \frac{r^4}{q^2}\longrightarrow0.
\]
The first condition will be proved sufficient for compactness; the second
makes the added smooth-interface contribution vanish. The final construction
chooses, for every $t\ge1$,
\[
 q=2^{16t},\quad m=2^{8t},\quad K=2^{6t},\quad
 R=2^{14t},\quad r=2^{7t}.
\]
Both $mq/r^4$ and $r^4/q^2$ equal $2^{-4t}$. The radius remains
$a=q^{-1/8}$ and the degree is exactly
\[
 D_U=|B_K\cap\mathbb Z^3|+|B_r\cap\mathbb Z^3|-1+6(t-1)
     =\Theta(q^{21/16}).
\]
The last term counts the remaining dyadic lengths between $r$ and $m$.
Over a horizon $1/c$, logical read incidences are $\Theta(q^{71/16})$,
compared with $\Theta(q^5)$ for the original dense radius-$\sqrt q$ family.
This fixes population, with slower radius convergence; it does not establish
equal-accuracy computational or native resource savings.

The threshold is necessary within this short-ball repair class. For the
periodic residue partition the added ball contributes exactly
\[
 X_{B_r}=\frac{2q^3}{m}C_r,
 \qquad C_r=\sum_{\|w\|\le r,\,w_1>0}w_1\sim\frac{\pi r^4}{4}.
\]
Its normalized entropy is asymptotic to $2r^4/(mq)$. If this ratio has
a bounded subsequence, the residue partitions have bounded repaired
entropy and no strong subsequence. The same conclusion follows after
localization inside $F$.

\paragraph{Compactness from the reads.}
Write $d(w)=\|u_h(\cdot+hw)-u_h\|_1$. For $\|s\|\le r/2$, average
$d(s)\le d(w)+d(s+w)$ over integer $\|w\|\le r/2$. This gives
\[
 d(s)\le Cr^{-3}\sum_{\|w\|\le r}d(w)
       \le Cqr^{-3}\mathcal E_{\rm fine}.
\]
Split translations of length $O(m)$ into $O(m/r)$ such steps. Fractional
cell translations are convex combinations of adjacent integer differences.
If $P_\epsilon$ is averaging over an $\epsilon$-cube, then
\[
 \|u_h-P_\epsilon u_h\|_1
 \le C\frac{mq}{r^4}\mathcal E_{\rm fine}.
\]
Choose a fixed smooth nonnegative bump supported in the half unit ball,
sample it on $v/K$, and normalize the weights $w_v$. They obey
$w_v\le CK^{-3}$ and $|w_v-w_{v-e_i}|\le CK^{-4}$. Put
$Qu=\sum_vw_vu(\cdot+\epsilon v)$ and $f_h=P_\epsilon Qu_h$.
Move the derivative of the box average onto the weight differences and
subtract the unshifted $u_h$, since those differences sum to zero. The
identity $mK^4=q^2$ yields
\[
 \|\nabla f_h\|_1\le C\mathcal E_{\rm coarse},
 \qquad
 \|u_h-f_h\|_1\le C\left(a+\frac{mq}{r^4}\right)\mathcal E_U.
\]
Each component energy is bounded by $\mathcal E_U$ separately; overlaps
are not charged twice. BV compactness therefore gives strongly convergent
subsequences for every bounded-entropy family, with binary finite-perimeter
limits. The explicit parameter choice proves this conclusion without
assuming compactness or filtering for successful interfaces.

\paragraph{Perimeter, exact volume and minimum shapes.}
For $u_h\to\mathbf1_E$ strongly, assign each $y\in B_1$ to
$v=\operatorname{round}(Ky)$. The distributional difference quotient and
lower semicontinuity give
\[
 \liminf a^{-1}\|u_h(\cdot+\epsilon v)-u_h\|_1
 \ge |D_y\mathbf1_E|.
\]
No convergence rate relative to $a$ is required: the difference quotient
acts on a fixed smooth test function after changing variables. Fatou and
the polar decomposition of the BV derivative imply
\[
 \liminf\mathcal E_U(u_h)
 \ge\frac2\pi\int_{B_1}|D_y\mathbf1_E|\,dy
 =\operatorname{Per}(E).
\]
For recovery, apply the directional BV translation formula to any set of
finite perimeter. It gives the coarse energy limit. Displacement-cell errors are
$O(K^{-1})\operatorname{Per}(E)$; the short ball adds at most
$C(r^4/q^2)\operatorname{Per}(E)$. Majority occupancy of each $h$-cell
has $L^1$ error $O(h)$ by the cube BV Poincare inequality, including the
fixed interior boundary restriction. Its energy error is
$O(D_U/q^2)\to0$. This gives binary recovery and the full perimeter
$\Gamma$-limit in strong $L^1$.

For prescribed $0<v<|C|$, the recovery needs only $O(q^2)$ label changes
to have exactly $\operatorname{round}(vq^3)$ occupied sites. A changed
label affects at most $D_U$ crossing pairs, so the additional normalized
cost is again $O(D_U/q^2)\to0$. Every finite problem has a minimum.
Compactness, the lower bound and recovery prove convergence of those
minimum values and of subsequences of all near-minimizers with $o(1)$
normalized excess. If a ball of volume $v$ fits strictly inside $C$, the
Euclidean isoperimetric equality case gives
\[
 \min\mathcal E_U\longrightarrow(36\pi)^{1/3}v^{2/3},
\]
and every limiting minimizer is a ball, up to translation and null sets.
Restoring units gives $(36\pi)^{1/3}V^{2/3}$ for physical volume $V$.
This identifies the limiting shape of the declared cut problem, not a
physical event horizon.

\paragraph{A uniform action and its hidden modes.}
The repair also permits equal positive action coefficients on every nonzero
read. Set $M_2=\sum_{v\in U}\|v\|^2$ and
\[
 A_U f=\alpha\sum_{v\in U\setminus\{0\}}[f(x)-f(x+hv)],
 \qquad \alpha=\frac6{h^2M_2}.
\]
Cubic symmetry gives exact continuum second moments and spatial error
at most $a^2B_4/4$, where $B_4$ bounds fourth directional derivatives of
the test function. No specially
weighted nearest-neighbor term is added. The uniform coefficient here
normalizes an action; the crossing entropy weights remain one.
On the clipped cube the zero-exterior convention includes a boundary
potential: if $b_x$ counts missing nonzero stencil offsets, then
\[
 \langle f,A_Uf\rangle_h=\alpha h^3\left[
  \sum_{\{x,y\}\,\mathrm{internal}}(f_x-f_y)^2+\sum_x b_x f_x^2\right].
\]
Its diagonal is $\alpha(|U|-1)$ at every site. The boundary term uses
prescribed zero exterior values and requires no outside record. Omitting it
gives a clipped graph Laplacian with a constant zero mode, rather than the
Dirichlet operator analyzed here. The finite evidence checks both the
constant and an affine field against an independent edge-energy formula.

For an integer ball of radius $\rho\ge8$ and principal phase $\theta$,
an elementary Dirichlet-sum estimate gives
\[
 \sum_{\|w\|\le\rho}[1-\cos(\theta\cdot w)]
 \ge\frac{\rho^3}{49152}\min(\rho^2\|\theta\|^2,1).
\]
To see uniformity away from zero, put an integer cube of side $N\ge\rho$
inside the ball. Its one-dimensional sum $d_N(x)$ satisfies
$1-|d_N(x)|/N\ge\min(N^2x^2,1)/16384$: for $|x|\le4/N$, expand
$N^2-|d_N|^2=2\sum_{d=1}^{N-1}(N-d)(1-\cos dx)$; beyond that interval
use $|d_N|/N\le\pi/(N|x|)$. Summing the other two coordinates proves
the displayed bound. Thus high-frequency control is not inferred from
selected numerical probes.

Write $s=r/K$ and decompose each Fourier index uniquely as
$k=m\ell+j$, $-m/2\le j_i<m/2$. The coarse phase depends only on $j$.
The moment bound $M_2\le56m^2K^5$ and the two ball-symbol bounds imply,
for $s\ge1$ and a fixed $C_0>0$,
\[
 \lambda_U(k)\ge C_0\min(\|j\|^2+s^5\|\ell\|^2,a^{-2}).
\]
The same conclusion holds when $s$ stays above any fixed positive constant,
with an adjusted bound. Since $K^4=m^3$, the entropy compactness parameter
is exactly $r^4/(mq)=s^4$.

If $s\to\gamma\in(0,\infty)$, second-moment expansion for every fixed
pair $(j,\ell)\in\mathbb Z^3\times\mathbb Z^3$ gives
\[
 \lambda_U(m\ell+j)\longrightarrow
 4\pi^2\bigl(\|j\|^2+\gamma^5\|\ell\|^2\bigr).
\]
The coarse integer phases cancel, the short-ball remainder vanishes because
$r/m\to0$, and the dyadic contribution is $O(n\alpha)\to0$.
This is the oscillator spectrum of a product of two three-tori with side
lengths $1$ and $\gamma^{-5/2}$. It describes additional spectral labels
of the regulator, not extra physical spacetime dimensions. Smooth fixed
spatial probes see only $\ell=0$.

For positive mass, the complete thermal limit follows as well. Below the
saturation scale the summands are bounded by a summable multiple of
$\exp[-c_0(\|j\|+\|\ell\|)]$; above it their total is at most
$Cq^3\exp(-c_0/a)\to0$. Thus the thermal partition function,
normal-ordered energy and Gibbs entropy converge to the product-torus
values at critical scale. If $s\to0$, arbitrarily many fixed aliases
approach the mass frequency and thermal energy diverges. If $s\to\infty$,
all $\ell\ne0$ contributions vanish under the same summable bound,
leaving the ordinary three-dimensional thermal limit. This last condition
is precisely the one that restores cut compactness.

For the explicit repair, $s^5=2^{5t}\ge a^{-2}=2^{4t}$, giving, for principal
physical wave vectors $p=2\pi k$, the stronger whole-spectrum estimate
\[
 C_1\min(\|p\|^2,a^{-2})\le\lambda_U(p)\le\|p\|^2.
\]
On the clipped cube, zero extension to a doubled periodic cube and this
bound give $L^2$ compactness and an $H_0^1$ limit. Positive-form duality,
consistency and smooth recovery give the full Dirichlet form limit;
min-max gives ordered eigenvalue convergence with multiplicities. The
zero extension has at most one eighth of its norm in the constant mode,
so there is a uniform Dirichlet gap, and compression yields
$\lambda_{n,j}\ge C_2\min(j^{2/3},a^{-2})$.
Summable low-frequency tails and the vanishing high-frequency total prove
the full Dirichlet thermal limits, including zero mass. Smooth-test action
norms control the momentum-covariance tail, giving the compact Gaussian
Weyl and Klein--Gordon commutator limits too.

All $T$ paths remain present. Charged coarse interpolation and binary
corrections give the local round causal limit. Use diamonds compactly
inside the clipped spatial cube and observation time slab; on the torus,
use non-wrapping diamonds, for example with time height $T$ satisfying
$cT<1/2$. The shrinking route tubes then stay inside the relevant patch.
Their generated-diamond volumes tend to $\pi c^3\tau^4/24$, where $\tau$
is the endpoint proper-time separation. The oriented strict-pair measure
is $V^2/20$. For $P$ causally ordered pairs among $N$ interval events,
$P/N^2\to1/20$ and the ordering fraction $2P/[N(N-1)]\to1/10$.
Fourth-root interval-count ratios give proper-time ratios.
Truncated and wrapping intervals do not obey these
unbounded flat-diamond formulas. The free
response cone agrees in the limit, including observations rounded to the
declared tick. Exact ODE evolution is not a native layer update: the
positive normalized trace forbids stable centered leapfrog at $a/c$.
Thermal Gibbs entropy and the crossing statistic remain distinct observables.

The complete proofs and explicit finite evidence are in
\ophid{code/m1_interfaces/DERIVATION.md} and
\ophid{code/m1_interfaces/receipt.json}. The independent checker reconstructs
all labels and unordered edges for 96 cut cases on 12 small graphs, covering
3,340,024 logical read incidences. Larger moment certificates use exact
column sums; their full populations are not claimed executed. Eight finite
graph and scale reductions have transitive standard-axiom audits. The
universal BV, isoperimetric and spectral arguments are analytic. The
continuum nonlocal-perimeter framework is standard; see A.~C.~Ponce,
\emph{A new approach to Sobolev spaces and connections to $\Gamma$-convergence},
Calculus of Variations and Partial Differential Equations 19 (2004), 229--255.
The atomic residue obstruction, its sharp sparse repair and the associated
critical spectral sector require the separate analysis given here.
